% Generated draft: requires mathematical and source review.
\documentclass[11pt]{article}
\usepackage[margin=25mm]{geometry}
\usepackage{fontspec}
\setmainfont{Latin Modern Roman}
\usepackage{amsmath,amssymb,mathtools,unicode-math}
\setmathfont{Latin Modern Math}
\usepackage{enumitem,needspace,xurl,hyperref}
\setlength{\parindent}{0pt}
\setlength{\parskip}{5pt}
\setlength{\emergencystretch}{3em}
\newcommand{\N}{\mathbb{N}}
\newcommand{\Z}{\mathbb{Z}}
\newcommand{\Q}{\mathbb{Q}}
\newcommand{\R}{\mathbb{R}}
\newcommand{\C}{\mathbb{C}}
\newcommand{\F}{\mathbb{F}}
\newcommand{\A}{\mathbb{A}}
\newcommand{\PP}{\mathbb{P}}
\newcommand{\E}{\mathbb{E}}
\newcommand{\eps}{\varepsilon}
\newcommand{\dd}{\,\mathrm{d}}
\begin{document}
\section*{UnsolvedMath ID: 11100009}
\textbf{Source problem code:} catalogue entry
\subsection{Definitions and mathematical statement}
Let \(\pi_{126}^S=\operatorname*{colim}_{j}\pi_{126+j}(S^j)\) be the stable homotopy group of spheres. By Pontryagin--Thom, its elements are represented by closed smooth stably framed \(126\)-manifolds. Their Kervaire invariant lies in \(\F_2\); in this dimension it is the Arf invariant of the framing-induced quadratic refinement of the middle-dimensional mod-two intersection form.

The question is whether there exists
\[\theta_6\in\pi_{126}^S,\qquad\operatorname{Kervaire}(\theta_6)=1.\]
Equivalently, in the mod-two Adams spectral sequence
\[E_2^{s,t}=\operatorname{Ext}_{\mathcal A}^{s,t}(\F_2,\F_2)\Longrightarrow(\pi_{t-s}^S)^\wedge_2,\]
where \(\mathcal A\) is the mod-two Steenrod algebra, does
\[h_6^2\in E_2^{2,128}\]
survive nontrivially to \(E_\infty\)? Here \(h_6\) has bidegree \((1,64)\), so the square has stem \(128-2=126\). Permanent survival means a nonzero permanent class, not merely the absence of one selected differential.

The supplied research record reports an affirmative result by Lin--Wang--Xu, \emph{On the last Kervaire invariant problem} (December 2024), proving this permanent survival and hence existence of \(\theta_6\). It says that together with the earlier Browder, Mahowald--Tangora, Barratt--Jones--Mahowald and Hill--Hopkins--Ravenel results, framed Kervaire-invariant-one manifolds exist exactly in dimensions \(2,6,14,30,62,126\). This is retained as the record's literature resolution, not a proof or independent verification here.
\subsection{Short English statement}
Does a framed Kervaire-invariant-one element exist in stable stem \(126\), equivalently survive as \(h_6^2\) in Adams filtration two? The record reports Lin--Wang--Xu's affirmative December 2024 result.
\subsection{Sources}
\begin{itemize}
\item Local UnsolvedMath registry, permanent ID \(11100009\), exact source code \texttt{catalogue entry}.
\item Mark Hovey, \emph{Hovey's Algebraic Topology Problem List}, Major problems, item 9, surviving mirror: \url{https://www-users.cse.umn.edu/~tlawson/hovey/}. Registry access 2026-07-29; html-ordered-list extraction.
\item W. Lin, G. Wang and Z. Xu, \emph{On the last Kervaire invariant problem}, arXiv:2412.10879 (December 2024): \url{https://arxiv.org/abs/2412.10879}, as supplied by the research record.
\item Browder, Mahowald--Tangora, Barratt--Jones--Mahowald and Hill--Hopkins--Ravenel are credited for the earlier cases/nonexistence results; no full entries are supplied.
\end{itemize}
\end{document}
